\section{Related Work}
\label{sec:related-work}
In this section, we review related literature on algorithm selection methods and their feature representations, graph representations of constraint problems, and graph feature extraction techniques.

\subsection{Algorithm Selectors and Handcrafted Features}
Most modern algorithm selectors employ machine learning techniques to predict which algorithm will perform best on an unseen problem instance. In the SAT domain, SATzilla \cite{xu2012satzilla2012} uses Random Forests over portfolios of SAT solvers. In CP, approaches like CPHydra \cite{bridge2012case} and SUNNY \cite{amadini2014sunny} utilise $k$-nearest neighbours to generate instance-specific solver schedules.

A central challenge in algorithm selection is constructing informative feature representations of problem instances. Traditionally, selectors have relied on manually designed, domain-specific features. For instance, SUNNY relies on \texttt{fzn2feat} \cite{Amadini2014-fn}, which extracts static features from a FlatZinc model, such as the counts and types of variables and constraints, domain sizes, and basic metrics derived from the variable-constraint graph (e.g., degree statistics). In its current version, \texttt{fzn2feat} comprises 95 features. However, such handcrafted features are expensive to design, often fail to generalize across distinct problem domains, and compress complex combinatorial substructures into flat summary statistics.

\subsection{Graph Representations for Constraint Problems}\label{sec:graph-representation}
While representing combinatorial problems as graphs is a conceptually intuitive approach \cite{bockmayr2005constraint}, determining the optimal structural configuration for such graphs remains a non-trivial challenge. Over the years, many different approaches have emerged: a lot of MIP-related tasks employ a bipartite graph representation where a pair of variable/constraint nodes share an edge when the variable has a non-zero coefficient in the constraint \cite{khalil2022mip}. Another approach, often used in CP, is to represent the problem as a tripartite graph, where nodes are constraints, variables and values in the variables' domains. Similarly to the bipartite version, variable and constraint nodes share edges if a variable has a non-zero coefficient in the constraint. Furthermore, variable nodes also have edges that connect them to all possible values in their domain \cite{cappart2023combinatorial}. Finally, \citet{boisvert2024towards} decide to decompose constraints into smaller components, and the final graph is akin to the abstract syntax tree of a program.

\subsection{Graph Kernels and GNNs in Combinatorial Optimisation}
Feature extraction transforms graph-structured data into fixed-dimensional vector spaces \cite{graphkernels}. Classic methods rely on graph-theoretical metrics to describe structural properties. Node-level features describe a node's topological importance, including Degree Centrality, Clustering Coefficient \cite{watts1998collective}, and Betweenness or Closeness Centrality \cite{newman2010network}. Edge-level features are primarily used for link prediction and node proximity, such as Common Neighbours, Jaccard Coefficient, Adamic-Adar Index \cite{adamic2003friends}, and Katz Index \cite{Luce_Perry_1949}.

Graph-level features are often extracted via graph kernels that compute similarity as a scalar product between graphs:
\begin{itemize}
    \item \textbf{Weisfeiler-Lehman (WL) Kernel} \cite{shervashidze2011weisfeiler}: Iteratively relabels nodes based on neighbor signatures (see Section \ref{sec:wl-feature_Extraction} for more details on the extraction process).
    \item \textbf{Graphlet Kernel} \cite{graphkernels}: Counts small, non-isomorphic subgraphs.
    \item \textbf{Random Walk Kernel} \cite{vishwanathan2010graph}: Compares graphs based on matching paths.
\end{itemize}
WL-based features are widely used in cheminformatics and computer vision \cite{young2022literature,simonovsky2017dynamic} as well as in neighbouring fields like automated planning \cite{chen2024return}. 

In recent years, Graph Neural Networks (GNNs) have gained significant popularity for combinatorial optimisation tasks. Crucially, the message-passing mechanism in standard GNNs is fundamentally connected to and bounded in expressive power by the 1-WL test \cite{Xuetal2018,morris2023weisfeiler}. However, while GNNs learn continuous representations via deep architectures, they require substantial training data and incur non-trivial inference latency when querying instances. For algorithm selection, where feature extraction and selector inference must be lightweight to avoid offsetting the performance gains of the selected solver, graph kernels based on Weisfeiler-Leman refinement provide a compelling alternative that preserves theoretical expressiveness while remaining fast to compute.
